%──────────────────────────────────────────────────────────────
% Gaurav Dhingra — Resume
% Compiler Engineer · LLVM · Systems Programming
%
% Build:  pdflatex resume.tex   (or just: make)
% Track:  git add resume.tex Makefile .gitignore
%──────────────────────────────────────────────────────────────
\documentclass[11pt, a4paper]{article}

% ── Packages ──
\usepackage[T1]{fontenc}
\usepackage[margin=0.6in, top=0.5in, bottom=0.5in]{geometry}
\usepackage{enumitem}
\usepackage{titlesec}
\usepackage[dvipsnames]{xcolor}
\usepackage{hyperref}
\usepackage{parskip}
\usepackage{tabularx}

% ── Colors ──
\definecolor{accent}{HTML}{2C5282}
\definecolor{muted}{HTML}{4A5568}
\definecolor{lightline}{HTML}{CBD5E0}

% ── Hyperlinks ──
\hypersetup{
    colorlinks=true,
    linkcolor=accent,
    urlcolor=accent,
    pdfauthor={Gaurav Dhingra},
    pdftitle={Gaurav Dhingra -- Resume},
}

% ── Section formatting ──
\titleformat{\section}
    {\large\bfseries\color{accent}}
    {}{0em}{}
    [\vspace{-0.6em}{\color{accent}\rule{\textwidth}{1.2pt}}]

\titlespacing*{\section}{0pt}{0.8em}{0.4em}

% ── List formatting ──
\setlist[itemize]{
    leftmargin=1.2em,
    itemsep=1pt,
    parsep=0pt,
    topsep=2pt,
    label=\textbullet,
}

% ── Custom commands ──
\newcommand{\role}[3]{%
    % #1 = Title — Org,  #2 = Location,  #3 = Dates
    \vspace{0.4em}
    \noindent\textbf{#1} \hfill \textcolor{muted}{\small #3}\\[-0.3em]
    \textit{\textcolor{muted}{\small #2}}
    \vspace{0.1em}
}

\newcommand{\ossproject}[1]{%
    \vspace{0.3em}\noindent\textbf{\small #1}
    \vspace{0.05em}
}

\newcommand{\skillrow}[2]{%
    \textbf{#1:} & #2 \\[2pt]
}

% ── No page numbers ──
\pagestyle{empty}

% ══════════════════════════════════════════════════════════════
\begin{document}

% ── Header ──
\begin{center}
    {\LARGE\bfseries Gaurav Dhingra}\\[0.2em]
    {\color{muted} Compiler Engineer\enspace·\enspace LLVM\enspace·\enspace Systems Programming}\\[0.4em]
    {\small
        \href{mailto:gauravdhingra.gxyd@gmail.com}{gauravdhingra.gxyd@gmail.com}
        \enspace|\enspace +91\,8791414504
        \enspace|\enspace \href{https://github.com/gxyd}{github.com/gxyd}
        \enspace|\enspace \href{https://gxyd.github.io}{gxyd.github.io}
    }
\end{center}

% ── Summary ──
\section{SUMMARY}

Compiler engineer with 8 years of software experience and deep expertise in
LLVM-based compiler infrastructure. Core contributor to LFortran (LLVM-based
Fortran compiler) and active contributor to LLVM/Clang. Background spans
compiler IR design, code generation, symbolic computation engines, and
expression parsing. Two-time Google Summer of Code participant with a sustained
track record of open-source contributions to SymPy, LFortran, and LLVM.

% ── Technical Skills ──
\section{TECHNICAL SKILLS}

\begin{tabularx}{\textwidth}{@{} l X @{}}
    \skillrow{Languages}{C++, Python, C, Kotlin, Bash}
    \skillrow{Compiler Infrastructure}{LLVM, Clang, clang-tidy, LLVM IR, ASR (Abstract Semantic Representation)}
    \skillrow{Tools \& Platforms}{Git, CMake, AddressSanitizer (ASan), GDB, Docker, CI/CD (GitHub Actions)}
    \skillrow{Domains}{Compiler optimization passes, frontend/backend code generation, symbolic computation, expression parsing, intermediate representations}
\end{tabularx}

% ── Experience ──
\section{EXPERIENCE}

\role{LLVM Project --- Open-Source Contributor}{Remote}{2025 -- Present}
\begin{itemize}
    \item Contributing to clang-tidy: developing and refining static analysis checks for the Clang tooling ecosystem
    \item Investigating and reproducing crash bugs in Clang's parser (e.g., stack-overflow in \texttt{ParsePostfixExpressionSuffix}) using locally built LLVM with AddressSanitizer on ARM64
    \item Engaging with the LLVM community through bug reports, code review, and patch submissions
\end{itemize}

\role{LFortran --- Compiler Engineer}{Remote}{March 2024 -- June 2025}
\begin{itemize}
    \item Core contributor to LFortran, a modern open-source Fortran compiler built on the LLVM backend, working toward its beta release
    \item Designed and improved the array operation pass in ASR (Abstract Semantic Representation), LFortran's high-level intermediate representation
    \item Implemented compiler frontend features and optimization passes in C++ using the LLVM C++ API
    \item Collaborated asynchronously with the open-source team on architecture decisions, code review, and CI workflows
\end{itemize}

\role{Digital Aristotle (Byju's) --- Senior Research Engineer}{Bangalore, India}{April 2022 -- February 2024}
\begin{itemize}
    \item Built a math solver engine for K-12 education: designed and implemented a parser and compiler pipeline that transformed LaTeX mathematical expressions into an AST and internal evaluation form
    \item Architected a domain-specific SDK in Kotlin for expression tree construction, type analysis, and step-by-step symbolic simplification
    \item Worked in a cross-functional agile team integrating the solver with GeoGebra's visualization layer
\end{itemize}

\role{Digital Aristotle (Byju's) --- Junior Research Engineer}{Bangalore, India}{June 2018 -- March 2022}
\begin{itemize}
    \item Developed the Autosolver backend for the Byju's BTLA app: a symbolic math engine using SymPy for automated step-by-step problem solving
    \item Designed and implemented a DSL for declarative step-by-step math solution templates, enabling content authors to define solver behavior without code changes
    \item Built CI/CD pipelines and test infrastructure in Python and TypeScript
\end{itemize}

% ── Open-Source Contributions ──
\section{OPEN-SOURCE CONTRIBUTIONS}

\ossproject{SymPy --- Pull Request Manager (NumFOCUS-funded)}
\begin{itemize}
    \item Selected as one of the top contributors to manage PR review and release workflows for SymPy (12k+ GitHub stars), a widely-used symbolic mathematics library. Position funded by NumFOCUS (Sep 2017, Dec 2017, Feb 2018)
\end{itemize}

\ossproject{Google Summer of Code 2017 --- SymPy}
\begin{itemize}
    \item Extended the Risch integration algorithm: implemented the parametric logarithmic derivative problem and added support for integrating trigonometric functions via the Risch algorithm
\end{itemize}

\ossproject{Google Summer of Code 2016 --- SymPy}
\begin{itemize}
    \item Implemented coset enumeration and Reidemeister--Schreier algorithms for finitely presented groups; added low-index subgroup computation for group-theoretic operations in SymPy
\end{itemize}

% ── Education ──
\section{EDUCATION}

\role{Indian Institute of Technology Roorkee}{}{2013 -- 2018}
Integrated M.Sc.\ in Applied Mathematics\\
{\small\textcolor{muted}{Relevant coursework: Design \& Analysis of Algorithms, Data Structures, Graph Theory, Linear Algebra, Discrete Mathematics}}

\end{document}
